\chapter{METODOLOGI PENELITIAN}
\section{Tahapan Penelitian}

\vspace{0.5cm}
\noindent\fbox{\begin{minipage}{\textwidth}
\textit{\textbf{Pedoman Penulisan:} Menjelaskan tahapan-tahapan yang dilakukan dalam melakukan penelitian ini.}
\end{minipage}}
\vspace{0.5cm}

\section{Instrumen Penelitian}

\vspace{0.5cm}
\noindent\fbox{\begin{minipage}{\textwidth}
\textit{\textbf{Pedoman Penulisan:} Menjelaskan alat bantu yang dipilih dan digunakan oleh peneliti dalam kegiatannya mengumpulkan data agar sistematis dan dipermudah.}
\end{minipage}}
\vspace{0.5cm}

\section{Metode Pengumpulan Data, Populasi, dan Sample Penelitian}

\vspace{0.5cm}
\noindent\fbox{\begin{minipage}{\textwidth}
\textit{\textbf{Pedoman Penulisan:} Menjelaskan metode pengumpulan data berupa pengamatan langsung, wawancara, studi pustaka, penentuan populasi, dan pengambilan sampel.}
\end{minipage}}
\vspace{0.5cm}

\section{Metode Analisis Data}

\vspace{0.5cm}
\noindent\fbox{\begin{minipage}{\textwidth}
\textit{\textbf{Pedoman Penulisan:} Menjelaskan tentang proses mengolah data sehingga menjadi informasi baru.}
\end{minipage}}
\vspace{0.5cm}

